\subsection{Breve estado del arte}

En la actualidad, la digitalización de la administración de propiedades ha cobrado impulso. Existen en el mercado plataformas comerciales de gran envergadura orientadas a corporaciones inmobiliarias. Sin embargo, estas herramientas de nivel empresarial suelen resultar económicamente inaccesibles para edificios pequeños o medianos, y presentan curvas de aprendizaje muy pronunciadas. En contraste, gran parte de las administraciones locales continúan operando de manera analógica, dependiendo de cuadernos físicos para el control de la guardia o de la distribución de expensas impresas. Esto evidencia una vacante para sistemas más ligeros, enfocados en resolver exclusivamente los problemas críticos diarios. CoConsorcio se posiciona en este segmento como una alternativa ágil, sirviendo como una prueba de concepto tecnológica ideal para aquellos consorcios que buscan dar su primer paso hacia la digitalización sin enfrentarse a sistemas sobredimensionados.
